% Limitations and Future Work: honest gaps and the roadmap.

The thesis of this paper depends on transparent VVUQ, so this section is
deliberately complete and does not soften any gap. It has two parts: an enumeration of what was approximated or could not be run, and a roadmap for
extending the assurance automation across oncology clinical trials.

\subsection{Limitations, approximations, and what could not be run}

The execution environment was a managed cloud container, and several inputs and
backends were absent by design or by circumstance. The agentic model backends
were not exercised live: the \texttt{anthropic} and \texttt{openai} packages
were not installed and no API key was present, so instruction and code generation
ran on their deterministic offline templates rather than a live model. Live web
search was not reachable, so ingestion exhausted its four attempt retry budget
and fell back to the clearly labeled offline stub, leaving the live HTML parsing
path unexercised. The PDF extraction backend was absent, so only the size based
chunk estimate ran and no real text extraction occurred. 

The physics and robotics
backends, namely the simulation engine, the NVIDIA Isaac stack, and ROS 2, were
not present and the container had no GPU or display, so the physical-ai modules
ran only as pure standard library references. The multibyte hard split chunking
edge case, which loses six characters at internal boundaries on a single
oversized line of multibyte text, was reported but not repaired, because the run
scope was to execute v0.1.0 rather than to modify the baseline. The LaTeX
template was not compiled to PDF in the container, since no TeX toolchain was
available there. 

The imagegen leg had no executable code in v0.1.0. And only a
single local Python version, 3.11.15, was run, with cross version behavior on
3.10 and 3.12 delegated to the continuous integration matrix. To these
environmental limits the discussion added two assurance gaps that matter more for
the thesis than any backend: there was no independent validation reference, so
the one accepted case was self referential, and the repository's rich input
corpus was not yet wired into the executed pipeline. These limitations are paired
with their remedies in \autoref{tab:limits_future}.

\begin{table}[H]
\caption{Limitations now and the tractable next steps.}
\label{tab:limits_future}
\centering
\begin{tabularx}{\textwidth}{L{6.2cm} >{\raggedright\arraybackslash}X}
\toprule
\textbf{Limitation (this run)} & \textbf{Future step} \\
\midrule
  No independent validation reference, so the accept is self referential & Supply an external gold standard so validation carries real weight. \\
  Input corpus not wired into the pipeline & Feed the chunked corpus through ingestion into codegen and papergen. \\
  Generation on offline templates, no live backends & Enable the live model, web search, PDF, and physics paths. \\
  Near deterministic 3 run simulation (CV 0.0007) & Genuine stochasticity, more runs, confidence intervals, sensitivity. \\
  Multibyte hard split loses 6 characters & Carry trailing bytes forward and add a multibyte chunking test. \\
  imagegen has no executable code in v0.1.0; template uncompiled & Close the codegen, imagegen, and papergen legs end to end. \\
  Single local Python version & Run 3.10, 3.11, 3.12 matrix locally, emit JSON results. \\
\bottomrule
\end{tabularx}
\end{table}

None of these limitations is hidden, and that is deliberate. The multibyte
chunking finding in particular surfaced only because the losslessness claim was
independently re verified rather than trusted, which is the same discipline the
gate applies to deliverables, now turned on the platform itself. A thesis that
asks the reader to value assurance over generation cannot be credible if it
conceals the places where assurance has not yet been applied, so each gap above
is paired in \autoref{tab:limits_future} with the concrete step that closes it,
and the most consequential gap, the absent independent validation reference, is
listed first because it bounds what any of the other results can claim.

\subsection{Future work: VVUQ automation across all oncology clinical trials}

The central forward claim is simple. This VVUQ automation should be adapted across
the full breadth of oncology clinical trials, because the speed, quality, and cost
benefits compound at scale while the gate preserves patient safety at every step.
A single trial gains a fast, auditable assurance loop; a portfolio of trials gains
a common standard of evidence that does not have to be rebuilt each time. The
tractable next steps follow in priority order, and they mirror the remedies in
\autoref{tab:limits_future} while adding the larger integration goals. Supply an
independent validation reference first, because without it the validation
dimension cannot do its job. Wire the existing input corpus through ingestion so
that code generation and paper generation operate on real source material.

Make the triple simulation genuinely stochastic, with more runs and reported
confidence intervals, so the dispersion bound is tested against real variation.
Enable the live model, network, PDF, and physics paths so the guarded fallbacks
become the exception rather than the rule. Repair the multibyte chunker and add a
multibyte test so the one known correctness gap is closed and stays closed. Close
all three legs of the thesis, code generation, image generation, and paper
generation, end to end, so the whole pipeline runs under one gate. 

Emit
machine readable results alongside the prose so the assurance decisions are
themselves auditable by downstream tools. The destination is a flexible, future
proof assurance layer that lets the industry check these steps off once rather
than solving them repeatedly, and it is designed to stay consistent as trials,
models, and procedures evolve~\cite{repo-cancer-automated, sel2025vvuq}. This is
the same direction as the author's federated and protocol to prediction oncology
work~\cite{kawchak2026federated, kawchak2025qsppdac}, and it is aligned with the
move toward faster, real time clinical evidence~\cite{fda2026realtime}.
